%% kapa: exam=1 round=20 year=2009 period=2 form=A subject=accounting subject_ko=회계학 seq=37 number=37
%% source: Q-Net 감정평가사 기출문제 (cst00309) artlSeq=5024178 file=extracted/1008/5024178/2009년 제20회 감정평가사 1차시험 기출문제/2교시.hwp
%% conversion: kordoc 4.24.51 -> tools/gen1.py (auto); review=proofread-sonnet; flags=none
\begin{question}{37}
(주)한국은 A와 B 두 개의 사업부를 이익중심점으로 설정하여 운영하고 있다. A 사업부는 부품을 생산하여 B 사업부와 기업외부에 판매할 수 있다. 사업부간의 대체가격은 단위당 변동제조원가에 기회원가를 가산하여 결정된다.
\begin{framedtext}
\begin{center}<A 사업부의 생산․판매 자료>\end{center}
\begin{tabular}{@{}l r@{\hspace{2em}}}
ㆍ연간 최대조업도 & 11,000단위 \\
ㆍ연간 고정제조원가 & ￦4,500,000 \\
ㆍ단위당 변동제조원가 & ￦1,800 \\
ㆍ단위당 외부시장 판매가격 & ￦3,000 \\
\end{tabular}
\end{framedtext}
A 사업부가 연간 9,000단위를 생산하여 전량 기업외부에 판매하고 있는 상황에서 B 사업부가 연간 4,000단위의 부품을 대체해 줄 것을 요청하였다. 다음 중 A 사업부가 요구해야 할 최소한의 단위당 대체가격은 얼마인가? (단, 대체거래를 하더라도 A사업부가 생산하는 부품의 제조원가는 주어진 자료와 동일하며 판매비와관리비는 고려하지 않는다. 또한 B사업부가 동일한 부품을 외부에서 구입하는 경우에는 단위당 ￦3,000을 지급하고 4,000단위 전량을 구입해야 한다.) \par
\begin{choices}
  \item ￦2,150
  \item ￦2,300
  \item ￦2,400
  \item ￦2,700
  \item ￦2,750
\end{choices}
\end{question}
\answer{3}
